% =========================================================================
% PoleSight — paper outline / index for future writing.
% Scaffold only (WACV has no rigid section layout); expand each bullet into
% prose later. Section flow mirrors RoadSight (ACM MM '26) dataset paper.
% =========================================================================
% \section*{Outline / Index}
\section{Introduction}
\label{sec:intro}
\begin{itemize}
    \item \textbf{Why poles.} Pole-like objects (light, power, traffic, gantry,
        fence) are key roadside/utility infrastructure for mapping, asset
        inventory, and AV localization.
    \item \textbf{How sensed.} Captured by mobile laser scanners as LiDAR point
        clouds alongside camera images.
    \item \textbf{Range images work.} Projecting LiDAR to 2D range images lets
        mature 2D detectors handle pole detection efficiently.
    \item \textbf{Bottleneck.} Collecting + labeling real scans is slow and
        costly $\rightarrow$ synthetic data helps scale.
    \item \textbf{Simulation is costly.} Full LiDAR simulation needs both a 3D
        environment \emph{and} a sensor-simulation step.
    \item \textbf{Our idea.} Simple synthetic range-image generation
        (project $\rightarrow$ augment), no simulator $\rightarrow$ scalable datasets.
\end{itemize}

\paragraph{Contributions.}
\begin{itemize}
    \item \textbf{PoleSight} dataset: multi-modal range images, five pole
        classes, $\sim$2000 labeled frames.
    \item Scalable \textbf{LiDAR$\rightarrow$range-image} synthetic generation method.
    \item Baselines: range-image detection on public data + diverse LiDAR setups.
\end{itemize}

\section{Related Work}
\label{sec:related}
\begin{itemize}
    \item \textbf{Domain augmentation \& simulated LiDAR.} % TODO cite
    \item \textbf{Range-image learning.} RangeNet++, YOLO-family detectors. % TODO cite
    \item \textbf{Urban point clouds \& mobile LiDAR.} Pole extraction from MLS. % TODO cite
    \item \textbf{Pole / infrastructure datasets.} Position PoleSight vs.\ prior. % TODO cite
\end{itemize}

\section{Method}
\label{sec:method}

\subsection{Range-Image Generation}
\begin{itemize}
    \item Initial: per-frame range-image generation + labeling.
    \item Augmentation loop: projection $\rightarrow$ generation $\rightarrow$
        validation $\rightarrow$ repeat.
\end{itemize}

\subsection{Detection Model}
\begin{itemize}
    \item Architecture + training setup.
    \item Validation protocol.
\end{itemize}

\section{Dataset and Experiments}
\label{sec:dataset}

\subsection{Dataset}
\begin{itemize}
    \item $\sim$2000 labeled frames.
    \item Classes: fence, gantry-sign, light, power, traffic poles.
    \item Modalities: RGB, IR/intensity, LiDAR range.
\end{itemize}

\subsection{Results}
\begin{itemize}
    \item ML evaluation: detection metrics, per-class.
    \item Comparison vs.\ sequential point-cloud baselines.
\end{itemize}

\section{Discussion}
\label{sec:discussion}
\begin{itemize}
    \item \textbf{Strengths.} Works on public data; diverse LiDAR setups; beats
        sequential point-cloud processing.
    \item \textbf{Weaknesses.} Static objects only; computational cost.
    \item \textbf{Future.} Improved augmentation; mixed real + synthetic data.
\end{itemize}

\section{Conclusion}
\label{sec:conclusion}
% One paragraph: dataset + scalable synthetic range-image method, takeaway.
